
%%%%%%%%%%%%%%%%%%%%%%%%%%%%%%%%%%%%
% Recap/Agenda
%%%%%%%%%%%%%%%%%%%%%%%%%%%%%%%%%%%%
% TODO better formatting
\begin{frame}
    \frametitle{Module 5: Beyond Recipes: Frequentist and Bayesian Approaches to Statistical Modeling}
    \begin{itemize}
        \item \hl{Previously: }Frequentist modeling using the likelihood function
        \item \hl{This time: }Bayesian modeling with discrete distributions
        \item \hl{Reading: }Bayes' Rules Chapter 3
    \end{itemize}

\end{frame}
%%%%%%%%%%%%%%%%%%%%%%%%%%%%%%%%%%%%
% Learning objectives:
%%%%%%%%%%%%%%%%%%%%%%%%%%%%%%%%%%%%
\begin{frame}
    \frametitle{Learning Objectives}
    \begin{itemize}
        \item \textbf{M5, LO2: The Bayesian Perspective:} Explain how Bayesian modeling uses likelihoods, priors, and posteriors to perform inference. Explain how Bayesian analyses can be updated sequentially as new data arrives and how the results remain invariant to the order in which the data is observed.
        \item \textbf{M5, LO4: Bayesian Modeling 1:} Construct Bayesian models for discrete distributions and identify the likelihood, prior, and posterior, and use the probability rules from Module 2 to compute conditional and marginal probabilities for the models.
    \end{itemize}
\end{frame}
